\begin{proposition}[4.2.7]
\label{IV.2.4.2.7-fr}
Sean $k$ un cuerpo, $X$ un $k$-preesquema, $K$ una extensión de $k$ tal que
$X\otimes_k K$ sea localmente noetheriano, $\mathcal F$ un
$\mathcal O_X$-módulo cuasicoherente, $x'$ un punto de $X\otimes_k K$, y $x$
su imagen en $X$.
\begin{enumerate}
\item[\textup{(i)}] Sea $(Z_\lambda)$ la familia de los subpreesquemas reducidos de $X$
cuyos espacios subyacentes son los ciclos primos asociados a $\mathcal F$~;
entonces las componentes irreducibles $Z_{\lambda\mu}$ de los
$Z_\lambda\otimes_k K$ son los ciclos primos asociados a
$\mathcal F\otimes_k K$, y $Z_{\lambda\mu}$ domina a $Z_\lambda$~; además,
para que un $Z_{\lambda\mu}$ sea inmerso, es necesario y suficiente que
$Z_\lambda$ lo sea.
\item[\textup{(ii)}] Para que $x$ pertenezca a un ciclo primo asociado inmerso de
$\mathcal F$, es necesario y suficiente que $x'$ pertenezca a un ciclo primo
asociado inmerso de $\mathcal F\otimes_k K$~; para que $\mathcal F$ carezca de
ciclos primos asociados inmersos, es necesario y suficiente que lo mismo
ocurra con $\mathcal F\otimes_k K$.

\oldpage[IV-2]{57}

\item[\textup{(iii)}] Los índices $\lambda$ tales que $x\in Z_\lambda$ son los mismos
que aquellos para los cuales existe un índice $\mu$ tal que
$x'\in Z_{\lambda\mu}$. En particular, si $x'$ no pertenece más que a un solo
ciclo primo asociado a $\mathcal F\otimes_k K$, $x$ no pertenece más que a un
solo ciclo primo asociado a $\mathcal F$.
\item[\textup{(iv)}] Si $X$ es localmente de tipo finito sobre $k$, se tiene
$\dim(Z_{\lambda\mu})=\dim(Z_\lambda)$.
\end{enumerate}
\end{proposition}

\begin{proof}
Se observará que la hipótesis implica que $X$ mismo es localmente noetheriano
(2.2.13 y 2.2.14)~; la afirmación (i) resulta de (4.2.5) y de la demostración
de (4.2.3), para $\mathcal G=\mathcal O_Y$, con
$Y=\operatorname{Spec}(K)$~; (ii) y (iii) resultan de (i) y de (2.3.5)~;
finalmente, (iv) es un caso particular de (4.2.4, (iii)).
\end{proof}

\begin{corollary}[4.2.8]
\label{IV.2.4.2.8-fr}
Si $X$ es localmente de tipo finito sobre $k$, el conjunto de las dimensiones
de los ciclos primos asociados es el mismo para $\mathcal F$ y
$\mathcal F\otimes_k K$~; el conjunto de las dimensiones de las componentes
irreducibles es el mismo para $X$ y $X\otimes_k K$.
\end{corollary}

\begin{proposition}[4.2.9]
\label{IV.2.4.2.9-fr}
Supongamos satisfechas las hipótesis de \textup{(4.2.5)}, y supongamos además que
$\mathcal F$ y $\mathcal G$ son coherentes. Sean
$(\mathcal F_\lambda)_{\lambda\in L}$ y
$(\mathcal G_\mu)_{\mu\in M}$ descomposiciones irredundantes de
$\mathcal F$ y $\mathcal G$, respectivamente~; para todo par
$(\lambda,\mu)\in L\times M$, sea
$(\mathcal H_{\lambda\mu\nu})_{\nu\in S(\lambda,\mu)}$ una descomposición
irredundante reducida de $\mathcal F_\lambda\otimes_k\mathcal G_\mu$, donde
$S(\lambda,\mu)=\operatorname{Ass}(\mathcal F_\lambda\otimes_k\mathcal G_\mu)$
\textup{(3.2.5)}. Entonces $(\mathcal H_{\lambda\mu\nu})$ \textup{(}para todos los triples
$(\lambda,\mu,\nu)$\textup{)} es una descomposición irredundante de
$\mathcal F\otimes_k\mathcal G$~; es reducida si
$(\mathcal F_\lambda)$ y $(\mathcal G_\mu)$ lo son.
\end{proposition}

\begin{proof}
Observemos que $\mathcal F\otimes_k\mathcal G$ y cada uno de los
$\mathcal F_\lambda\otimes_k\mathcal G_\mu$ son coherentes, y cada uno de los
$\mathcal F_\lambda\otimes_k\mathcal G_\mu$ se identifica con un cociente de
$\mathcal F\otimes_k\mathcal G$~; cada uno de los
$\mathcal H_{\lambda\mu\nu}$ se identifica por definición con un cociente
coherente de $\mathcal F_\lambda\otimes_k\mathcal G_\mu$, y por tanto también
con un cociente coherente de $\mathcal F\otimes_k\mathcal G$. La familia de
los soportes de los $\mathcal H_{\lambda\mu\nu}$ es localmente finita, pues
$\operatorname{Supp}(\mathcal F_\lambda\otimes_k\mathcal G_\mu)$ está contenido
en el espacio subyacente a
$\operatorname{Supp}(\mathcal F_\lambda)\times_k
\operatorname{Supp}(\mathcal G_\mu)$ (donde
$\operatorname{Supp}(\mathcal F_\lambda)$ y
$\operatorname{Supp}(\mathcal G_\mu)$ designan los subpreesquemas cerrados
reducidos correspondientes), y para $\lambda$ y $\mu$ dados, la familia de los
soportes de los $\mathcal H_{\lambda\mu\nu}$ es localmente finita por
hipótesis~; nuestra afirmación resulta, por tanto, de que las familias
$(\operatorname{Supp}(\mathcal F_\lambda))$ y
$(\operatorname{Supp}(\mathcal G_\mu))$ son localmente finitas. Para demostrar
que $(\mathcal H_{\lambda\mu\nu})$ es una descomposición irredundante de
$\mathcal F\otimes_k\mathcal G$, basta, por tanto, demostrar que el
homomorfismo canónico
$\mathcal F\otimes_k\mathcal G
\longrightarrow\bigoplus_{\lambda,\mu,\nu}\mathcal H_{\lambda\mu\nu}$
es inyectivo~; ahora bien, es el compuesto de los homomorfismos
$\mathcal F\otimes_k\mathcal G
\longrightarrow\bigoplus_{\lambda,\mu}(\mathcal F_\lambda\otimes_k\mathcal G_\mu)
\longrightarrow\bigoplus_{\lambda,\mu,\nu}\mathcal H_{\lambda\mu\nu}$~;
el segundo de estos homomorfismos es inyectivo por definición, y lo mismo
ocurre con el primero, que no es sino el producto tensorial de los
homomorfismos canónicos inyectivos
$\mathcal F\to\bigoplus_\lambda\mathcal F_\lambda$,
$\mathcal G\to\bigoplus_\mu\mathcal G_\mu$ (recuérdese que
$\mathcal F$ y $\mathcal G$ son \emph{planos} sobre $k$). Finalmente, si
$(\mathcal F_\lambda)$ y $(\mathcal G_\mu)$ son reducidas, se puede suponer
que $L=\operatorname{Ass}(\mathcal F)$ y
$M=\operatorname{Ass}(\mathcal G)$~; el hecho de que los
$\mathcal H_{\lambda\mu\nu}$ formen entonces una descomposición reducida
resulta de que los $\operatorname{Ass}(\mathcal F_\lambda\otimes_k
\mathcal G_\mu)$ forman una partición de
$\operatorname{Ass}(\mathcal F\otimes\mathcal G)$ en virtud de (4.2.2)
(cf. (3.2.5)).
\end{proof}

\begin{corollary}[4.2.10]
\label{IV.2.4.2.10-fr}
Bajo las hipótesis de \textup{(4.2.7)}, supongamos además que $\mathcal F$ es coherente,
y sea $(\mathcal F_\lambda)_{\lambda\in L}$ una descomposición irredundante de
$\mathcal F$~; para todo $\lambda\in L$, sea
$(\mathcal F_{\lambda\mu})_{\mu\in
\operatorname{Ass}(\mathcal F_\lambda\otimes_k K)}$ una descomposición
irredundante reducida de $\mathcal F_\lambda\otimes_k K$. Entonces
$(\mathcal F_{\lambda\mu})$ es una descomposición irredundante de
$\mathcal F\otimes_k K$~; es reducida si $(\mathcal F_\lambda)$ lo es.
\end{corollary}

\begin{proof}
Se aplica (4.2.9) con $Y=\operatorname{Spec}(K)$,
$\mathcal G=\mathcal O_Y=K$, y $M$ reducido a un solo elemento.
\end{proof}

\oldpage[IV-2]{58}

